\subsection*{Answer selection under independently authored wording}
An independent local model wrote 96 Chinese requests about 24 held-out headwords. Each call specified the requested fact and historical stage. The author saw neither training questions nor answer images. A writer receiving broader wording supervision recovered 66 complete images, up from 21 for its predecessor. It answered all four requests correctly for seven headwords, up from zero (Fig.~\ref{fig:independent-questions}). It answered 37 of 48 direct questions and 29 of 48 questions commissioned with known facts or forms as distractors. Removing four requests judged ambiguous before scoring gave 62 of 92 complete answers. The local author is a synthetic source of phrasing.

With wider supervision, the shared-update writer also recovered 4,540 of the 4,544 familiar requests with the status before its field, up from 356. This gain did not extend to two reserved constructions. Of 1,136 distinct requests that explicitly corrected a previous choice, four yielded complete correct answers. Of 4,544 requests referring to a field and form as the former or the latter, 1,134 were correct. No group of four anaphoric requests was entirely correct. The candidate-specific update gave 52 of 96 independently authored answers and three complete headwords despite ending with a lower training loss. Neither writer passed the registered generalization criterion. In a subsequent control, repeating the original questions for the same 5,576 updates and target order gave 18 of 96 exact answers. Thus the expanded wording gave 48 additional exact answers over matched replay. Removing the greeting from the illustrated failure restored the requested headword but still elicited its pronunciation in place of its original meaning.
